% !TeX program = pdflatex
% !TeX encoding = UTF-8
% pxLST code inline & annex test: \pxCodeInline, \printPxCodeAnnex only (no list commands).
% Compile from tex/ with: TEXINPUTS=.;./latex;./code;./tests pdflatex -interaction=nonstopmode tex/tests/pxLST.list-code.test.tex
% Or run via: powershell -File tests/run-pxLST-tests.ps1 (from tex/).
\documentclass{article}
\usepackage{pxLST}
\begin{document}

\section*{pxLST code inline \& annex test}

\subsection*{1. Inline code (1--2 snippets, trailing empty slots)}
Code inline: \pxCodeInline{int x = 0;}{return x + 1;}{}{}{}{}{}{}.

\subsection*{2. Opt-in: no print commands here}
This subsection uses only inline \pxCodeInline{baz}{}{}{}{}{}{}{} and never calls \printPxCodeAnnex. No annex output; no errors.

\subsection*{3. Named annex-id (code scopes)}
\pxCodeInline[annex-id=scope-a]{snippet A}{}{}{}{}{}{}{}
\pxCodeInline[annex-id=scope-b]{snippet B}{}{}{}{}{}{}{}

\section*{Annexes (default scope)}
\printPxCodeAnnex[title=Code snippets (default)]

\section*{Annexes (named scopes)}
\printPxCodeAnnex[annex-id=scope-a, title=Annex A code]
\printPxCodeAnnex[annex-id=scope-b, title=Annex B code]

\end{document}
